% Pursuing Stacks (cotes 134-2 to 134-8): archive cote 134-7, site page 96 (PDF page 97), typescript page 538.
% typescript. Draft compared: ps6.tex:4826-4910.
% AI-transcribed by Claude Opus 5.5 (claude-opus-5-5, adaptive thinking, effort high), Message Batches API, batch ps4 (msgbatch_012Aoc1NuskATPsyBeKSYK7J), 2026-11-12; reviewed against both supplied scan halves; see CHECK_DECISIONS.tsv.
\scanpage{134-7/96}
\mspage{538}

homotopy sheaves $\underline{\hand{\Pi}}_i$ satisfy the relations
\[ (6) \qquad H^i(T, \underline{\hand{\Pi}}_j(X_*)) = 0 \ \text{ for } \ i > 0 \quad , \]
as a consequence maybe of \xout{a}\inserted[typed]{some} spectral sequence whose abu\over[typed]{\unsure{m}}{t}ment is the graded
homotopy of the cointegration of $X_*$ over T. If so, then (5) are quasi-%
i\unsure{s|d}omorphisms, whenever \xout{we} the sheaves $\underline{\hand{\Pi}}_i$ can be endowed with the struc-%
ture of a quasi-coherent $\underline{O}_k$-Module.

\hand{128}. Next requirement about the notion of "bundles" we are going to work
with, is about existence of a ``linearization functor'' ; associating to
every bundle another bundle, namely it's ``linearization'' L(X), endowed
moreover with the structure of an $\underline{O}_k$-Module. \uline{We insist that for a given}
\uline{bundle,} L(X) \uline{should be at any rate a bundle too}. More specifically,
if U(k) denotes the category of all bundles over k, we should also int\over[typed]{\illegible{1}}{ro}-%
duce a corresponding ``k-linear'' category, or more accurately a $\underline{O}_k$-linear
category, Uab(k), whose objects should be objects of U(k) endowed
with the extra structure of an $\underline{O}_k$-Module, possibly subject to some \xout{restri}
restrictions - \xout{thus we} thus we'll get a forgetful functor
\[ (7) \qquad \hand{K} \ ; \ \mathrm{Uab}(k) \hand{\longrightarrow} U(k) \quad . \]
When we take $\over[typed]{\unsure{B}}{U}(k)$ to be unipotent bundles in the sense of section 111,
the evident choice for Uab(k) is to take the category of quasicoherent
$\underline{O}_k$-Modules, equivalent to the category of k-modules $\mathrm{Ab}_k$ . Dually,
we may take U(k) to be the category of sheaves over k isomorphic to the
underlying sheaf of sets of a vector bundle \xout{over k} V(M) over k associated
to a k-module M , by the requirement
\[ (8) \qquad V(M)(k') = \mathrm{Hom}_{k\text{-mod}}(M,k') \quad . \]
T\over[typed]{h}{h}e evident corresponding choice for Uab(k) is to take the category of
all vector bundles over k, which is equivalent to the category $\mathrm{Ab}_k^{o}$
\uline{opposite} to the category of k-modules, as the functor $M \hand{\mapsto} V(M)$ is contra-%
variant in M. Maybe we should distinguish between these two choices
of bundles by different notations, namely
\[ (9) \qquad U_W(k) \ \text{ and } \ U_V(\overset{\unsure{\hand{?}}}{k}) \quad , \]
where the subscripts W and V are meant to suggest the standard descripti-%
ons of objects via \hand{(}the \deleted{(}underlying sheaves of sets of the $\underline{O}_k$-Modules)
\xout{W(M) and V(M) respectively}
\[ W(M) \ \text{ and } \ V(M) \]
respectively, where \hand{(we recall)}
\[ (10) \qquad W(M)(k') = M \hand{\otimes}_k k' \ . \]
\margin{96}

% Visual collation: both supplied image halves inspected; see CHECK_DECISIONS_134_7_8.tsv.
